% Parts A--B. Working contribution; all result IDs in this file are owned by AB.
\section{A. The group-theoretic assertions of section 1}


\begin{result}{A1}{Known}{The exact cohomological dimensions and the real exception}
\label{res:A1}
Let $K/\mathbb Q$ be finitely generated of transcendence degree $n$, let
$k$ be the algebraic closure of $\mathbb Q$ in $K$, and put
$P=G_{K\overline k}$ and $E=G_K$. With continuous cohomology of discrete
torsion modules,
\[
 \operatorname{cd}_\ell P=n\quad\hbox{for every prime }\ell,
 \qquad \operatorname{cd}P=n,
\]
\[
 \operatorname{cd}_\ell E=n+2\quad(\ell\ne2),\qquad
 \operatorname{cd}_2E=
 \begin{cases}n+2&K\text{ not formally real},\\
 \infty&K\text{ formally real}.
 \end{cases}
\]
In every case $\operatorname{vcd}_2E=\operatorname{cd}_2G_{K(i)}=n+2$.
Thus the equalities proposed in F1 hold, with precisely this exception
at two (149, section 1, manuscript p.~6, scan 13). For a general
function field, the relevant condition is an ordering of $K$ itself,
not just a real embedding of its constant number field.
\end{result}
\begin{proof}
The opened upper-bound theorem is
\cite[Theorem 15.2]{AB_MilneLEC}: a field of transcendence degree $d$
over $F$ has cohomological dimension at most $\operatorname{cd}F+d$.
For number fields the precise input is
\cite[Theorem I.4.10(c)]{AB_MilneADT}: in degrees at least three,
cohomology with finite coefficients is the product of the cohomology
at real places. Thus a totally imaginary number field has dimension
at most two. Apply these bounds to $K\overline k/\overline k$ and
to $K(i)/k(i)$ to obtain $n$ and $n+2$, respectively.

We give the lower bound and the comparison to $K$. On a smooth
$n$-dimensional model over $k$, choose a closed point and a complete
regular flag through it. The successive orders along the flag define
a valuation of rank $n$, value group $\mathbb Z^n$ with lexicographic
order, and residue field a finite extension $k'/k$. The same construction
over $\overline k$ has residue field $\overline k$. Its henselization
is algebraic over the original function field, hence its absolute group
embeds as a closed subgroup of that field's absolute group.

At each rank-one stage, a henselian discrete valuation with residue
characteristic zero has inertia $\widehat{\mathbb Z}(1)$: finite
extensions of the strict henselization are obtained by roots of a
uniformizer, since units have all roots by Hensel's lemma. The
Hochschild--Serre sequence therefore has only rows zero and one.
The Kummer class of a uniformizer restricts to the generator of the
inertia row. Cup product with that class splits the residue map
\[
 H^{j+1}(L,\mu_\ell^{\otimes(r+1)})
       \longrightarrow H^j(\kappa,\mu_\ell^{\otimes r})
       \qquad(r\in\mathbb Z).
\]
This proves in particular that a nonzero residue class remains nonzero
after successive cups by the $n$ uniformizers. Starting with $1$ in
$H^0(\overline k,\mathbb Z/\ell)$ gives degree $n$ in the geometric
case. Starting with a nonzero class in $H^2(k',\mu_\ell)$ gives a nonzero
class of degree $n+2$ with coefficients $\mu_\ell^{\otimes(n+1)}$
in the arithmetic case. Such a class exists: localization onto
$H^2(k'_v,\mu_\ell)\simeq\mathbb Z/\ell$ is surjective by
\cite[Theorem I.2.1 and Corollary I.4.16]{AB_MilneADT}, taking $v$
a finite place and leaving another place available. Monotonicity of
cohomological dimension for closed subgroups supplies both lower bounds.
This construction also handles $n=0$, with no uniformizers.

The exact finite-index theorem we opened is Serre's unnumbered
Theorem in section 1 of\cite{AB_Serre1965}: for a profinite group
without an element of order $\ell$, an open subgroup has the same
$\ell$-cohomological dimension. By Artin--Schreier (A2), $G_K$ has
no odd-prime torsion. It has no involution exactly when $K$ is not
formally real: an involution fixes a real closure, and any ordering
has a real closure whose absolute group has order two. Serre's theorem
applied to $G_{K(i)}\subset G_K$ proves the finite values above.
When $K$ is formally real, the closed subgroup of order two has
nonzero mod-two cohomology in arbitrarily high degrees, so
$\operatorname{cd}_2G_K=\infty$. This also proves the virtual formula.
\end{proof}

\subsection{Finite order, free groups, and surface groups}

\begin{result}{A2}{Known}{Finite-order elements and complex conjugation}
\label{res:A2}
For a field $F$ of characteristic zero, every nonidentity element of finite
order in $G_F=\operatorname{Gal}(F^{\rm sep}/F)$ has order two. Its centralizer
is the subgroup of order two that it generates. In $G_{\mathbb Q}$ all such
elements are conjugate to the complex conjugation determined by
$\overline{\mathbb Q}\subset\mathbb C$. Every open subgroup of
$G_{\mathbb Q}$ has trivial centralizer in $G_{\mathbb Q}$.
This answers F2 (149, section 1, manuscript p.~6, scan 14).
\end{result}
\begin{proof}
The exact Artin--Schreier input is the finite-absolute-group theorem:
a nontrivial finite absolute Galois group has order two; its fixed field is
real closed, of characteristic zero, and its algebraic closure is obtained by
adjoining $\sqrt{-1}$. We opened this statement in
\cite[section 3, Example 2, p.~22]{AB_Geyer2012}.
For $z\in G_F$ of finite order, the fixed field of $\langle z\rangle$ has
absolute group $\langle z\rangle$, so the theorem applies.

Put $R=(F^{\rm sep})^{\langle z\rangle}$. An element commuting with $z$
preserves $R$ and restricts to an $F$-automorphism of $R$. The order on $R$
is determined by its squares. This automorphism therefore preserves the
order. For each polynomial in $F[T]$, its finitely many roots in $R$ form
an ordered finite set, and an order-preserving permutation of this set is
the identity. Since $R/F$ is algebraic, the automorphism fixes every element
of $R$. Thus it belongs to $\operatorname{Gal}(R(\sqrt{-1})/R)
=\langle z\rangle$. The reverse inclusion is immediate.

The ordered field $\mathbb Q$ has its unique usual order. Its algebraic real
closures are uniquely isomorphic as ordered extensions: a root is mapped
to the corresponding ordered root of its minimal polynomial. Extending
such an embedding to algebraic closures shows that the subgroups of order
two in $G_{\mathbb Q}$ are conjugate.

For the last assertion we use the opened
\cite[Theorem 2.7 and Proposition 1.2]{AB_MinamideTsujimura2022}:
a nontrivial closed normal subgroup of the absolute Galois group of a
Hilbertian field has trivial centralizer. We use the author version dated
15 February 2022, in which the theorem number is 2.7. We also checked
Statement B2 of the supplied return\cite{AB_Cyclotomic}; its normal
centralizer argument agrees with this exact primary theorem.
The normal core $N$ of an open $H\subset G_{\mathbb Q}$ is open and
nontrivial. The field $\mathbb Q$ is Hilbertian
\cite[section 4, Theorem 1]{AB_Geyer2012}. Hence
$C_{G_{\mathbb Q}}(H)\subset C_{G_{\mathbb Q}}(N)=1$.
In particular $Z(H)=1$, including the open normal subgroups mentioned on
scan 14. This establishes every input in the manuscript's finite-centralizer
argument, without assuming its conclusion.
\end{proof}

\begin{result}{A3}{Proved}{Open subgroups, torsion, and centers of profinite surface groups}
\label{res:A3}
Let $P$ be the full profinite completion of either a free group of finite
rank $r\geq2$, or the fundamental group of a closed oriented surface of
genus $g\geq2$. Every open subgroup has trivial centralizer in $P$, and
$P$ has no nonidentity element of finite order. An open subgroup of index
$d$ is the full profinite completion of a free group of rank
$1+d(r-1)$, respectively a surface group of genus $1+d(g-1)$.
The free rank-one exception is $\widehat{\mathbb Z}$, whose center and the
centralizers of all its open subgroups equal the whole group.
These statements answer F3 (149, section 1, manuscript pp.~7 and 13,
scans 15 and 28).
\end{result}
\begin{proof}
Write $P=\widehat\Lambda$. For an open subgroup $H$, its inverse image
$\Lambda_H\subset\Lambda$ has index $d$. Its full profinite topology is
the induced one: for a finite-index subgroup $V\subset\Lambda_H$, the
intersection of its finitely many conjugates in $\Lambda$ is a finite-index
subgroup of $\Lambda$ contained in $V$. Consequently
$\widehat{\Lambda_H}\to H$ is an isomorphism. The finite covering of the
rose with $r$ loops, respectively of the oriented surface, has Euler
characteristic multiplied by $d$. A spanning tree in the graph gives the
stated free basis; the classification of compact oriented surfaces gives
the stated genus. Its abelianization is therefore a finite product of
$\widehat{\mathbb Z}$ and is torsion-free.

This proves torsion-freeness of $P$ itself by a finite-quotient argument.
For a nontrivial finite-order $z\in P$, choose open normal $N$ with
$z\notin N$ and put $V=N\langle z\rangle$. The image of $z$ in
$V^{\rm ab}$ must be zero, because this abelianization is torsion-free.
But $V\to V/N$ is a nontrivial cyclic quotient generated by $zN$, so it
factors through $V^{\rm ab}$ and gives a contradiction.

We also give the center argument explicitly. Let $N\triangleleft P$ be
open and $D=P/N$. The corresponding regular covering has cellular chain
modules $\mathbb Q[D]^r\to\mathbb Q[D]$ in the graph case, and
$\mathbb Q[D]\to\mathbb Q[D]^{2g}\to\mathbb Q[D]$ in the closed surface
case. The $H_0$ representation is trivial in both cases; $H_2$ is also
trivial in the oriented surface case. Additivity of characters in these
finite complexes gives
\[
 H_1\simeq\mathbb Q\oplus\mathbb Q[D]^{r-1},\qquad
 H_1\simeq\mathbb Q^2\oplus\mathbb Q[D]^{2g-2},
\]
respectively, as rational $D$-representations. These identities follow
from the character equalities and semisimplicity of finite-group
representations in characteristic zero. The representations are faithful:
a nonidentity $d\in D$ has trace $1$, respectively $2$, whereas the
identity has strictly larger trace. Comparison with the maximal pro-$\ell$ quotient
$N^{\rm ab,(\ell)}\otimes_{\mathbb Z_\ell}\mathbb Q_\ell$ gives the same action for any prime
$\ell$. A central $z\in P$ acts trivially on $N^{\rm ab}$, hence has
trivial image in $D$. This holds for every $N$, so $z=1$.

Apply the same argument to any open normal $N\subset H$. The kernel of
$C_P(N)\to P/N$ is $Z(N)=1$, so the centralizer is finite. Torsion-freeness
makes it trivial, and $C_P(H)\subset C_P(N)$ finishes the proof.
The literature statements with exactly these full profinite hypotheses
are Nakamura's Corollary (1.3.4)\cite{AB_Nakamura1994}, and
Mochizuki--Tamagawa's Proposition 1.4 and Remarks 1.1.5, 1.2.2
\cite{AB_MochizukiTamagawa2008}; we opened their statements. The above
argument supplies the center, open-subgroup, and torsion reductions used
here.
For completeness, the nonhyperbolic geometric curve types in characteristic
zero have groups $1$ for $(0,0)$ and $(0,1)$,
$\widehat{\mathbb Z}(1)$ for $(0,2)$, and
$\widehat{\mathbb Z}^{\,2}$ for $(1,0)$. The latter two have their entire
group as the centralizer of every open subgroup. The local meridian is
trivial for $(0,1)$; in $(0,2)$ the two meridians generate the same group,
with opposite orientations. These statements follow directly from the
sphere, punctured-sphere, and torus presentations.
\end{proof}

\begin{result}{A4}{Proved}{Inertia at distinct lifted cusps}
\label{res:A4}
Let $C$ be a smooth connected curve over an algebraically closed field of
characteristic zero, of hyperbolic type $(g,r)$ with $r>0$, and let
$P=\pi_1(C)$. In the normalization of its smooth compactification in the
universal pro-\'etale cover, each cusp $c$ has inertia
$I_c\simeq\widehat{\mathbb Z}(1)$. Distinct lifted cusps have
$I_c\cap I_d=1$. Moreover
\[
 N_P(I_c)=I_c,\qquad C_P(H)=I_c\quad(1\ne H\subset I_c).
\]
Here subgroups are actual chosen subgroups; changing the universal lift
conjugates all statements simultaneously. This supplies the loop inputs
requested with F3.
\end{result}
\begin{proof}
For a complex curve the local punctured disk gives $I_c=\widehat{\mathbb Z}(1)$;
finite covers and algebraically closed base change transport the statement
to the stated fields. To see injectivity without discarding the one-cusp
case, pass to a finite unramified cover of the compact curve when $g>0$:
the quotient of its surface group by a nonzero map to $\mathbb Z/3$ gives
at least three preimages of every cusp. For $g=0$ hyperbolicity already
gives at least three cusps. A cusp in a curve with at least two punctures
maps injectively into its abelianized geometric group: the puncture part is
\[
 \widehat{\mathbb Z}(1)^{S}/
 \widehat{\mathbb Z}(1)(1,\ldots,1).
\]
This formula follows by abelianizing the usual oriented surface
presentation; the only relation among meridians is their sum.

For distinct pro-cusps $c,d$, a finite cover separates their images.
Refine it, as just described, until the compactification has at least three
punctures. If $1\ne h\in I_c\cap I_d$, a positive integer power of $h$
lies in the open subgroup of this covering and is still nontrivial, by
torsion-freeness of $\widehat{\mathbb Z}$. Its two images in the displayed
abelianization belong to different coordinate lines. Those lines have
zero intersection for $|S|\geq3$: comparison of a third coordinate in a
multiple of the diagonal forces that multiple, and then both coordinates,
to be zero. Injectivity of the inertia line gives a contradiction.

An element normalizing $I_c$, or centralizing its nontrivial subgroup,
therefore fixes the lifted cusp. Its stabilizer is $I_c$, since the residue
field at a geometric cusp is algebraically closed. Conversely $I_c$ is
abelian. This proves both equalities. For comparison we opened
\cite[Lemma 1.3.7]{AB_Mochizuki2004}, which states commensurable
terminality for both the full profinite group and each pro-$\ell$ group;
the argument here also establishes the stronger intersection assertion
used in this paper.
\end{proof}

\begin{result}{A5}{Proved}{Extensions and inverse limits}
\label{res:A5}
The property $C_P(H)=1$ for every open subgroup $H\subset P$ passes to
extensions of profinite groups and to directed inverse limits with
surjective transition maps. Consequently every successive extension of
nonabelian free profinite groups and profinite hyperbolic surface groups
has this property. The limit statement fills the omission between scans
12 and 15 (149, section 1, manuscript pp.~5--7).
\end{result}
\begin{proof}
For $1\to A\to B\xrightarrow q D\to1$, let $H\subset B$ be open.
An element centralizing $H$ has image centralizing the open subgroup
$q(H)$ of $D$, so belongs to $A$. It then centralizes $H\cap A$, open
in $A$, so is the identity.

For $P=\varprojlim P_i$, the projections $q_i$ are continuous surjections
of compact groups and are open. Thus $q_i(H)$ is open for every open $H$.
Every coordinate of an element of $C_P(H)$ belongs to
$C_{P_i}(q_i(H))=1$, so the element is the identity.
For the manuscript's models the surjectivity is concrete: a connected
finite \'etale cover of a normal integral scheme stays connected on a
dense open, since the function field of the cover is unchanged. Hence
$\pi_1(V)\to\pi_1(U)$ is surjective for dense opens $V\subset U$.
Every finite separable extension of the function field is finite \'etale
over a sufficiently small open of its normalization. Finite-cover
correspondence therefore gives
\[
 G_{k(U)}\simeq\varprojlim_{V\subset U}\pi_1(V).
\]
The identical construction after algebraic closure of the constant field
gives the geometric comparison. These are isomorphisms of profinite
groups, not merely bijections of abstract groups.
\end{proof}

\begin{result}{A6}{Known}{Theorem 1 and its strengthening}
\label{res:A6}
Let $K/\mathbb Q$ be finitely generated, let $k$ be the algebraic closure
of $\mathbb Q$ in $K$, and fix an algebraic closure $\overline K$ with
$\overline k\subset\overline K$. Put
$E=G_K$ and $P=G_{K\overline k}$. Every open subgroup of $E$, respectively
$P$, has trivial centralizer in $E$, respectively $P$. For
$\operatorname{trdeg}(K/k)>0$ one has, more strongly,
\[
 C_E(H)=1\quad\text{for every open }H\subset P,
 \qquad G_k\longrightarrow\operatorname{Out}(P)\text{ is injective}.
\]
This is Theorem 1 and the Remarque (149, section 1, manuscript pp.~6 and
10, scans 13 and 21; F4).
\end{result}
\begin{proof}
The Hilbertian normal-centralizer theorem used in A2, independently
verified in\cite[Theorem 2.7 and Proposition 1.2]{AB_MinamideTsujimura2022}
and checked against the attached return\cite{AB_Cyclotomic}, applies to $K$: it is an infinite finitely generated
field and hence Hilbertian. It also applies to $K\overline k$ when the
transcendence degree is positive, because this field is finite over a
rational function field of positive transcendence degree over
$\overline k$. The required Hilbertian facts are
\cite[section 4, Theorem 2(a),(d),(e)]{AB_Geyer2012}. Applying that theorem to the
normal core of each open subgroup proves the two first assertions. In
degree zero $P=1$ and its assertion is tautological.

Here is the reduction to a curve, including its comparison maps. Choose
$t\in K$ transcendental over $k$, and let $F$ be the relative algebraic
closure of $k(t)$ in $K$. The extension $F/k(t)$ is finite, $k$ is
algebraically closed in $F$, and $K/F$ is regular. Let $C/k$ be the smooth
projective geometrically connected curve with function field $F$, and
choose a nonempty affine hyperbolic open $U\subset C$. Restriction on
separable closures and then restriction to covers unramified on $U$ give
continuous surjections
\[
 P\twoheadrightarrow G_{F\overline k}
   \twoheadrightarrow\pi_1(U_{\overline k}).
\]
The first is surjective because regularity makes $K$ linearly disjoint
from every finite separable extension of $F$; the second follows from
finite-cover correspondence. Both maps are induced over $k$, so commute
with the outer $G_k$-actions. An inner action on $P$ descends to an inner
action on this quotient. The opened Hoshi--Mochizuki
\cite[Theorem C(ii)]{AB_HoshiMochizuki2011} says that, for a hyperbolic
curve over a perfect characteristic-zero field, the kernel of the full
outer representation is contained in the kernel of the outer map to
$G_{\mathbb Q}$; in particular it is zero over a number field.
Here $k$ is a number field. This proves injectivity of $G_k\to\operatorname{Out}(P)$.
Since $Z(P)=1$, the kernel of this action identifies with $C_E(P)$:
an inner action of a lift $e$ can be corrected by its unique inducing
element of $P$ to centralize $P$.

Finally let $H\subset P$ be open. There is an open normal subgroup
$E_0\subset E$ with $N=E_0\cap P\subset H$: this follows from the
subspace topology on the closed subgroup $P$. The group $N$ is a
nontrivial closed normal subgroup of $E_0=G_{K_0}$, where $K_0/K$ is
finite. It is nontrivial because $P$ is infinite and $N$ has finite index
in $P$. The same theorem gives $C_{E_0}(N)=1$, hence
$C_E(H)\cap E_0=1$. Therefore $C_E(H)$ embeds in the finite group
$E/E_0$. By A2 a nonidentity element of this centralizer would be an
involution whose centralizer is its own order-two group. It cannot
centralize the infinite group $H$. Thus $C_E(H)=1$.
This last argument supplies the step needed after shrinking the
arithmetic group in the manuscript; membership in the shrunken group
was not silently assumed.
\end{proof}

\section{B. Rigidity, the translation error, and Kummer theory}

\begin{result}{B1}{Proved}{The rational-point group of a semiabelian variety}
\label{res:B1}
Let $K$ be finitely generated over its prime field and let
$1\to T\to G\to A\to1$ be a semiabelian variety over $K$.
As an abstract abelian group $G(K)$ is a direct sum of a finite group and
a free abelian group, whose rank can be infinite. Consequently
\[
 \bigcap_{n\geq1}nG(K)=0.
\]
In characteristic $p>0$ the exact prime-to-$p$ variant is
\[
 \bigcap_{(n,p)=1}nG(K)=G(K)[p^\infty].
\]
The finite-generation assertion on scan 34 is false: already
$\mathbb G_m(\mathbb Q)=\mathbb Q^*$ has one independent generator for
each rational prime. This answers the erroneous input of F6 (149,
section 1, manuscript p.~16, scan 34).
\end{result}
\begin{proof}
Let $L/K$ be a finite splitting field of $T$. We first prove that
$L^*$ is the direct sum of its finite roots of unity and a free abelian
group. Write $k_0$ for the algebraic closure of the prime field in $L$.
A normal projective geometrically integral model $Z/k_0$ has
\[
 1\longrightarrow k_0^*\longrightarrow L^*
 \xrightarrow{\operatorname{div}}\operatorname{Div}(Z).
\]
The kernel consists of global invertible regular functions on the proper
integral model, hence of constants. The divisor group is free abelian;
its subgroup $\operatorname{im}(\operatorname{div})$ is free, so this
sequence splits as a sequence of abstract abelian groups onto its image.
For finite $k_0$, its multiplicative group is finite. For a number field
$k_0$, the analogous divisor sequence over its ring of integers, together
with the opened unit theorem\cite[Theorem 5.1]{AB_MilneANT}, expresses $k_0^*$ as finite torsion plus a free
abelian group. Thus the same holds for $L^*$.

The injection $T(K)\hookrightarrow T(L)\simeq(L^*)^{\dim T}$ now proves
the same assertion for $T(K)$: its torsion is finite and its quotient by
torsion is a subgroup of a free abelian group, hence free. The exact
Mordell--Weil--Lang--N\'eron statement needed for the quotient is
\cite[Theorem 16.7]{AB_MilneAV}: an abelian variety over a field finitely
generated over the prime field has a finitely generated rational-point
group. Hence the image $B$ of $G(K)\to A(K)$ is finitely generated.

Choose finitely many lifts in $G(K)$ of generators of $B$, and let $H$
be the subgroup they generate. The finitely generated subgroup
$H\cap T(K)$ is contained in a finitely generated direct summand $T_0$
of $T(K)$ which also contains its torsion: in a free basis only finitely
many coordinates occur. Write $T(K)=T_0\oplus T_1$. Then
\[
 G(K)=(H+T_0)\oplus T_1.
\]
Indeed every element is a sum from these groups, and an element of their
intersection belongs to $H\cap T(K)+T_0\subset T_0$ and to $T_1$, so
vanishes. The first summand is finitely generated, and the second is
free. The assertions about intersections of multiples follow coordinate
by coordinate. For a finite abelian group, multiplication by integers
prime to $p$ is invertible on its $p$-primary part and eventually kills
every other primary part. The rational-prime valuations on $\mathbb Q^*$
verify the advertised counterexample to finite generation.
\end{proof}

\begin{result}{B2}{Proved}{Kummer comparison and Theorem 3 in characteristic zero}
\label{res:B2}
Let $K/\mathbb Q$ be finitely generated, let $X/K$ be geometrically
connected of finite type, and let $i:X\hookrightarrow G$ be an immersion
into a semiabelian variety. Both maps (*) and (**) on scans 32--33 are
injective: $K$-points have distinct geometric conjugacy classes of
sections, already after abelianizing the geometric kernel. This is
Theorem 3 in its stated embedding generality (149, section 1, manuscript
pp.~15--17, scans 32--35; F6).
\end{result}
\begin{proof}
For $n\geq1$, multiplication by $n$ on $G$ is a finite \'etale torsor
under $G[n]$. Pullback along $x\in G(K)$ gives its Kummer class
\[
 \delta_n(x)\in H^1(G_K,G[n]),\qquad
 \ker\delta_n=nG(K).
\]
The geometric monodromy of this torsor is abelian. Therefore its
pullback by a point section depends only on the section in
$E_G/\overline{[\pi_1(G_{\overline K}),\pi_1(G_{\overline K})]}$.
For two point sections this finite-level comparison is exactly
$\delta_n(x-y)$: with division points $n\widetilde x=x$ and
$n\widetilde y=y$, the two cocycles differ by
$\sigma(\widetilde x-\widetilde y)-(
\widetilde x-\widetilde y)$.
Consequently equality of the abelianized section classes puts $x-y$
in every $nG(K)$, and B1 gives $x=y$.

For completeness the converse comparison, over any characteristic-zero
field, also holds. The natural map
$\pi_1(G_{\overline K})\to T(G)=\varprojlim_nG[n]$ is an isomorphism:
after descent to a finitely generated subfield and embedding in $\mathbb C$,
a complex semiabelian group has fundamental group $\mathbb Z^{2\dim A+\dim T}$,
and its multiplication coverings form a cofinal system of finite connected
coverings. Algebraically closed base change and the finite-cover comparison
transport the statement back. If $x-y\in\bigcap_n nG(K)$, the finite
nonempty sets of $K$-rational division points, indexed by divisibility,
have a compatible point. To check this without a surjectivity assumption,
every finite family of compatibility conditions is solved at one common
multiple; compactness of the product of the finite sets then solves all
conditions. Thus the full Tate-module Kummer cocycle is a coboundary.
This rechecks the equivalence stated as D7 in the attached return.

Applying the morphism of arithmetic extensions induced by $i$ sends equal
section classes for $X$ to equal classes for $G$. Equality of $i(x),i(y)$
implies $x=y$ because an immersion is a monomorphism. Passage from full
sections to their geometric abelianizations is the displayed quotient
homomorphism, so injectivity of (**) proves that of (*).

Every prime is needed for a theorem uniform in $K,G$. On omitting a
prime $q$, take $K=\mathbb Q(\zeta_q)$ and $G=\mathbb G_m$.
The points $1$ and $\zeta_q$ differ by $q$-torsion, which is divisible by
every integer prime to $q$. Their prime-to-$q$ Kummer classes coincide.
More exactly, B1 shows that the kernel on retaining a set $\Sigma$ of
primes is the torsion of $G(K)$ supported outside $\Sigma$. Thus the
manuscript's Remarque is a uniform necessity statement; it does not say
that every individual torsion-free example requires every prime.
\end{proof}


\begin{result}{B3}{Proved}{Corollary 1(a), with the translation in the correct group}
\label{res:B3}
Let $k$ be algebraically closed, let $X$ be a reduced connected
finite-type $k$-scheme, and let $f,g:X\to Y$ be morphisms, where
$i:Y\hookrightarrow G$ is an immersion into a semiabelian variety.
For any prime $\ell\ne\operatorname{char}k$, equality
$H_1(f,\mathbb Z_\ell)=H_1(g,\mathbb Z_\ell)$ implies a unique
$a\in G(k)$ with $i\circ g=\tau_a\circ i\circ f$.
This proves Corollary 1(a), including its one-prime refinement, on
scans 24--25 (149, section 1, manuscript pp.~11--12; F5).
The printed/transcribed $a\in Y$ must be corrected to $a\in G$:
with $Y=\mathbb G_m\setminus\{1\}$ and $f=g$, the difference is
$1\notin Y$.
\end{result}
\begin{proof}
Put $h=i\circ g-i\circ f$. The addition law on a commutative algebraic
group induces addition on its abelianized prime-to-characteristic
fundamental groups: this follows by pulling back the $[\ell^n]$
Kummer torsors, whose cocycles add. Hence $H_1(h,\mathbb Z_\ell)=0$.
We show that such a map $h:X\to G$ is constant.

It suffices first to consider a smooth connected curve $U$ mapping
to $X$. The composite $U\to G\to A$ extends to the smooth proper
compactification $C$ by properness of $A$. After subtracting its
value at a chosen point, it factors through a homomorphism
$J_C\to A$: the exact opened universal property is
\cite[III, Proposition 6.1]{AB_MilneAV}. The map
$H_1(U,\mathbb Z_\ell)\to H_1(C,\mathbb Z_\ell)$ is surjective,
since a connected finite \'etale cover of $C$ stays connected on a
dense open. The Abel map identifies the latter module with
$T_\ell J_C$\cite[III, Proposition 9.1 and Corollary 9.6]{AB_MilneAV}.
Thus the homomorphism $J_C\to A$ is zero on $T_\ell J_C$ and is zero.
For this last implication, a positive-dimensional image would have a
kernel of smaller dimension; a finite union of translates of that
kernel contains at most a constant times $\ell^{2nd}$ points of
order dividing $\ell^n$, with $d<\dim J_C$, and cannot contain all
$\ell^{2n\dim J_C}$ such points for every $n$.

The map $U\to G$ therefore lies in one torus fibre. Translate it into
$T\simeq\mathbb G_m^r$. The inclusion of $T$ induces an injection
on $\ell$-adic homology: the exact sequences on torsion give
$0\to T_\ell T\to T_\ell G\to T_\ell A\to0$ (surjectivity uses
divisibility of $T(k)$). Each coordinate is a unit on $U$, and its
map on a meridian at a boundary point is multiplication by its
integer order there, as is seen from the torsor $z^{\ell^n}=h$.
Zero homology therefore forces every boundary order to vanish.
The coordinate extends to an invertible function on proper $C$ and
is constant. This proves the curve assertion.

A nonconstant morphism on an irreducible finite-type component is
nonconstant on some integral curve: take an affine open, choose a
nonconstant coordinate from the image, and complete it to a
transcendence basis; after localization, the finite map to affine
space has a curve lying over a line on which that coordinate varies.
Normalize a nonempty open of this curve to apply the preceding
argument. Thus $h$ is constant on every irreducible component of
$X$. Their intersection graph is connected, and constants on
intersecting components coincide. Finally, equality on all closed
points of a reduced finite-type scheme implies equality of the
morphisms, since the ideal of their closed equalizer has zero
radical and is zero. This proves constancy on $X$. Any $x\in X(k)$
then gives $a=h(x)$ and proves uniqueness.
\end{proof}

\begin{result}{B4}{Refuted}{The open-subset translation lemma and its exact curve correction}
\label{res:B4}
The lemma of scan 37 is false for the hyperbolic type $(1,1)$.
For a smooth hyperbolic curve $Y=C\setminus D$ over an algebraically
closed field, with its generalized Abel map $\alpha:Y\hookrightarrow J_D^1$,
the following is the exact correction on curve types: a translation
sending a nonempty open of $\alpha(Y)$ into $\alpha(Y)$ is trivial for
every hyperbolic type except $(1,1)$; for type $(1,1)$ every translation
has that property. Full setwise stability $\alpha(Y)+u=\alpha(Y)$ forces
$u=0$ in every hyperbolic type. Here the field has characteristic zero;
the same proof uses prime-to-characteristic cohomology in arbitrary
characteristic. This answers the lemma and the homological proof in F5
(149, section 1, manuscript pp.~18--19, scans 37--40).
\end{result}
\begin{proof}
Let $E:y^2=x^3-x$, with origin $O$ at infinity, and put
$Y=E\setminus\{O\}$. Its generalized Jacobian for the reduced one-point
modulus is $E$ itself. Choose $u=(0,0)$, a nonzero point of order two,
and $U=E\setminus\{O,u\}$. Translation by $u$ interchanges $O,u$ and
therefore preserves $U$, so $U+u\subset Y$ although $u\ne0$.
On $U$ it is explicitly
\[
 (x,y)\longmapsto(-x^{-1},yx^{-2}).
\]
Substitution verifies the equation of $E$ and applying this map twice
gives the identity. The script \texttt{code/AB\_checks.py} checks these
identities symbolically and checks the point permutation over $\mathbb F_5$.
In general, for $Y=E\setminus\{c\}$ and any $u\in E$, the open
$E\setminus\{c,c-u\}$ proves the asserted property of every translation.

We prove the positive statement and identify the failed step. Translation
on the generalized Jacobian induces the identity on $H_1(-,\mathbb Z_\ell)$
for $\ell$ invertible in the field. The generalized Abel map identifies
$H_1(Y,\mathbb Z_\ell)$ with $T_\ell(J_D)$; this follows from the
degree-zero divisor construction, or from the compact Jacobian Tate
module and the meridian exact sequence
\[
 0\longrightarrow\mathbb Z_\ell(1)^D/\mathbb Z_\ell(1)
 \longrightarrow H_1(Y,\mathbb Z_\ell)
 \longrightarrow T_\ell(J_C)\longrightarrow0.
\]
The translation on an open part therefore extends to an automorphism
$f:C\xrightarrow\sim C$ whose action on $T_\ell(J_C)$ is the identity.
The extension is birational because the translated map has inverse given
by the opposite translation on its image; a birational map between smooth
proper curves extends uniquely in both directions.

For $g(C)\geq2$, the group of automorphisms of $C$ is finite: its
faithful pluricanonical projective representation is an algebraic group
with tangent space $H^0(C,T_C)=0$, since $\deg T_C=2-2g<0$.
Trivial action on a Tate module implies triviality of the induced
Jacobian homomorphism by the torsion-count argument in B3. Its action
on the Jacobian is faithful: if an automorphism acts trivially there,
its Abel image is translated by a torsion point $a$. The cyclic group
generated by a nontrivial such translation acts freely on $C$. Its quotient
is \'etale of degree $m>1$, while invariance of $H_1$ makes the quotient
have the same genus, by the averaging projector on rational $\ell$-adic
cohomology. The unramified formula
$2g(C)-2=m(2g(C/\langle f\rangle)-2)$ contradicts $g(C)>1$.
Thus $f=1$ and the immersion gives $u=0$.

For the remaining hyperbolic types assume $r=|D|\geq2$ (and $r\geq3$
in genus zero). Shrink $U$ so $J=C\setminus U$ contains both $D$ and
$f^{-1}D$. For the two maps $U\to Y$, the kernels on the meridian
module $M_J=\mathbb Z_\ell(1)^J/\mathbb Z_\ell(1)$ are generated by
the coordinates in $J\setminus D$ and $J\setminus f^{-1}D$.
Equality of the homology maps makes these submodules equal. For subsets
of cardinal at most $|J|-2$, their coordinate spans in $M_J$ are equal
only when the subsets themselves are equal: choose a coordinate outside
their union when available, or compare the coordinates in the two
complements directly; a nonzero multiple of the diagonal has all
coordinates nonzero and cannot express one missing coordinate using
fewer than $|J|-1$ others. Hence $f(D)=D$.
It follows that $f$ acts trivially on $H_1(Y)$ and therefore on $M_D$.
The permutation representation on $M_D$ is faithful for $|D|\geq2$;
for two points the transposition is multiplication by $-1$, which is
different from $1$ on $\mathbb Z_\ell$. Thus $f$ fixes each point of $D$.
In genus zero three fixed points force $f=1$. In genus one its trivial
Jacobian action makes it a translation, and a translation with a fixed
point is trivial. Again $u=0$ follows.

When $|D|=1$, both killed subsets have $|J|-1$ elements, and each of
their spans is all of $M_J$. This is precisely the exceptional case
recognized algebraically on scan 39 but not eliminated on scan 40:
shrinking $U$ does not change the number $|D|=1$. For the explicit elliptic
example the two maps even induce the same map on $H_1$, whereas the full
fundamental-group maps differ. A meridian about $u$ is killed by the
inclusion $U\to Y$ and maps to the nontrivial meridian about $O$ under
translation. A conjugation cannot turn the identity into a nonidentity.
This locates the exact nonabelian information lost in the manuscript's
homological argument.

For full stability in type $(1,1)$, the translation must preserve the
unique missing point, so it is trivial. This completes the classification.
\end{proof}


\begin{result}{B5}{Proved}{Corollaries 1(b), 2(c)--(d), and 3 with all comparison maps}
\label{res:B5}
A hyperbolic polycurve means a finite tower of smooth pairs whose
geometric curve fibres have $2g-2+r>0$, including proper fibres.
Over a field finitely generated over $\mathbb Q$, rational points
of a hyperbolic polycurve have distinct geometric conjugacy classes
of full point sections, and a point section has trivial centralizer
in the geometric group. For a reduced geometrically connected
finite-type source, the arithmetic outer fundamental-group map
therefore determines a morphism into that polycurve. The same
morphism rigidity holds for a target immersed in a semiabelian
variety, already after geometrically abelianizing both extensions.

Over an algebraically closed field of characteristic zero, two maps
$f,g:X\to Y$ with equal geometric outer maps are equal in each of
these cases: $Y$ is a hyperbolic polycurve and $f$ has open image
on fundamental groups; or $Y$ is immersed in a semiabelian variety
and $C_{\pi_1(Y)}(f_*\pi_1(X))=1$. In particular a dominant map
from an irreducible reduced source into a hyperbolic polycurve is
rigid in this sense. These are Corollaries 1(b), 2(c)--(d), and 3
(149, section 1, manuscript pp.~11--15, scans 24--31; F5).
\end{result}
\begin{proof}
We first prove the centralizer statement for a curve $V/K$ and any
section $s:G_K\to E_V$. Write $\Delta=\pi_1(V_{\overline K})$.
For $1\ne z\in C_\Delta(s(G_K))$, choose a characteristic open
normal subgroup $N\subset\Delta$ excluding $z$; such subgroups are
cofinal because $\Delta$ is finitely generated. Put
$H=N\overline{\langle z\rangle}$. Both $N$ and $H$ are stable under
$s(G_K)$, so $Hs(G_K)$ defines a geometrically connected finite
\'etale curve cover $V_H/K$. The image of $z$ in $H^{\rm ab}$ is
nonzero, since its image in the finite cyclic quotient $H/N$ is
nonzero. But it is fixed by $G_K$. The abelian group $H^{\rm ab}$
is the full Tate module of the generalized Jacobian of $V_H$,
with its natural $G_K$-action. Its invariants are zero: a compatible
system of rational torsion points would be infinitely divisible
in the finite torsion group supplied by B1. This contradiction
proves the assertion. The identification with the Tate module is
the compact Jacobian comparison together with the meridian module;
it is compatible with the arithmetic action because inner
conjugation acts trivially on abelianization.

Point injectivity for a curve follows from B2 and its generalized
Abel immersion. Prove point injectivity and the point-section
centralizer assertion for a tower by induction. For its last map
$Y\to B$, C4 supplies the exact sequence and, for $b\in B(K)$,
the actual comparison isomorphism
\[
 E_{Y_b}\ \xrightarrow{\sim}\ E_Y\times_{E_B}s_b(G_K).
\]
It is the identity on the geometric fibre kernel and on $G_K$.
Equal point sections for $y,z\in Y(K)$ first imply equality of their
images $b\in B(K)$. The projection of a geometric conjugator then
centralizes $s_b(G_K)$, hence is one by induction. It lies in the
geometric fibre kernel; the displayed map reduces the equality to
point sections of the curve $Y_b$, where B2 applies. For a centralizer
element of $s_y$, its projection is likewise one, and the curve
centralizer assertion kills the remaining fibre element. This gives
both induction steps. No immersion of the entire tower was used.

For arithmetic morphism rigidity, restrict two equal extension maps
along a point section $s_x$ after each finite field extension $K'/K$.
The Cartesian field-base-change comparison
$E_{X_{K'}}=E_X\times_{G_K}G_{K'}$ makes their values the sections
of $f(x),g(x)$. Point injectivity gives equality at all such points.
These include all geometric closed points. Reducedness and
separatedness give equality of the maps, as in B3. For a semiabelian
immersed target the same proof uses B2 and works on abelianized
geometric kernels. If ``conjugation by $\pi_1(Y)$'' on scan 30 means
the entire arithmetic group, its arithmetic image centralizes $G_K$,
because both maps lie over $G_K$ and the source surjects onto $G_K$.
A6 gives $Z(G_K)=1$, so the conjugator is geometric after all.

Now descend the geometric data to a field $K$ finitely generated over
$\mathbb Q$ inside the given algebraically closed field. Algebraically
closed base change preserves these full geometric fundamental groups.
Use the given geometric conjugator to make the restrictions of the
two arithmetic maps $F,G:E_X\rightrightarrows E_Y$ equal on
$\Delta_X$. For $e\in E_X$, the element $G(e)F(e)^{-1}$ lies in
$\Delta_Y$ and centralizes $F(\Delta_X)$: conjugate an element
$F(\delta)=G(\delta)$ by both sides and use equality on
$e\delta e^{-1}\in\Delta_X$. Thus triviality of this centralizer
makes $F=G$. Arithmetic rigidity just proved gives $f=g$.
For a polycurve, A3, A5, and C4 give trivial centralizers of open
geometric subgroups, so this proves case (c) as well as (d).

Finally, a dominant morphism $f:X\to Y$ with $X$ integral and
$Y$ normal has open fundamental-group image. The restriction map
$G_{k(X)}\to G_{k(Y)}$ has image the open subgroup belonging to the
finite relative algebraic closure of $k(Y)$ in $k(X)$. The map
$G_{k(Y)}\to\pi_1(Y)$ is surjective, because a connected finite
\'etale cover of a normal integral scheme has connected generic
fibre. The resulting open subgroup of $\pi_1(Y)$ lies in
$f_*\pi_1(X)$ by functoriality of the generic point. A subgroup
containing an open subgroup is open. This proves Corollary 1(b).
The argument proves these geometric statements for reduced
connected sources in the stated cases; the normality assumption
on $X$ in Corollary 2 is unnecessary for this proof.
The tentative margin on scan 24 that removes every dominance or image
condition is false, as the author himself observes on scan 26.
For $Y=\mathbb P^1\setminus\{0,1,\infty\}$ and any nonempty reduced
connected source, the constant maps with values $2$ and $3$ are distinct
and both induce the trivial geometric fundamental-group homomorphism.
Their image centralizer is all of $\pi_1(Y)$, so B5 states the precise
additional hypotheses used by the surviving corollaries. This records
the withdrawn assertion as sweep row ZAB7 (manuscript pp.~11--12).
\end{proof}

\subsection{\texorpdfstring{Characteristic $p$: full and geometrically prime-to-$p$ groups}{Characteristic p: full and geometrically prime-to-p groups}}

\begin{result}{B6}{Proved}{The full point-section theorem in positive characteristic}
\label{res:B6}
Let $K/\mathbb F_p$ be finitely generated. If a geometrically connected
finite-type $K$-scheme $X$ admits an immersion into a semiabelian variety,
then its $K$-points give distinct geometric conjugacy classes of sections,
already after abelianizing the full geometric \'etale fundamental group.
Thus both (*) and (**) of Theorem 3 remain injective with the full groups.
This proves the full-group assertion and the finite-field specialization
margin on scan 35 (149, section 1, manuscript p.~17; F7).
\end{result}
\begin{proof}
First let $A/K$ be an abelian variety and let $a\in A(K)$ have the same
abelianized section class as $0$. Spread $A,a$ over a smooth integral
finite-type $\mathbb F_p$-scheme $S$ with function field $K$, shrinking
so that $A_S$ is an abelian scheme. If $a\ne0$, the complement of the
equalizer of $a,0$ is a nonempty open and contains a closed point $s$.
Its residue field $k_s=\mathbb F_q$ is finite and $a_s\ne0$. Set
$R=\mathcal O_{S,s}^{\rm h}$ and $L=\operatorname{Frac}R$.
The inclusion $K\hookrightarrow L$ induces restriction of section classes.

The map $F_q-1:A_s\to A_s$ is a finite \'etale surjection with constant
kernel $H=A_s(\mathbb F_q)$. Its derivative is $-1$, hence it is \'etale;
properness makes it finite, and its image is nonempty, open and closed
in the connected variety $A_s$. Its geometric kernel consists of the
$q$-Frobenius-fixed points. Its fibre above $0$ is a trivial $H$-torsor,
whereas its fibre above $a_s\ne0$ has no $\mathbb F_q$-point, because
$F_q(z)-z=0$ for every such point $z$.

The precise opened proper henselian comparison is this: for a proper
scheme over a henselian local ring, restriction to the closed fibre is
an equivalence on finite \'etale covers
\cite[Lemma 58.9.1, Tag 0A48]{AB_Stacks}.
It lifts the Lang cover to a finite \'etale $H_R$-torsor $V\to A_R$;
full faithfulness lifts the action and torsor identity. The pullback to
the zero section is trivial, by the henselian equivalence on finite
\'etale $R$-schemes\cite[Lemma 58.8.2, Tag 09ZS]{AB_Stacks}.
This torsor has abelian geometric monodromy. Equality of the
geometrically abelianized point-section classes therefore makes its
pullbacks by $0,a\in A(L)$ isomorphic. Hence $a^*V$ has an $L$-point.
A finite \'etale scheme over the normal domain $R$ with an $L$-point
has an $R$-point: its coordinate values are integral over $R$ and in
$L$, and thus belong to $R$. Reduction gives an $\mathbb F_q$-point
of the Lang fibre above $a_s$, a contradiction. Thus $a=0$.

For $1\to T\to G\to A\to1$, equal abelianized point sections for
$x,y\in G(K)$ give
$d=x-y\in\bigcap_{(n,p)=1}nG(K)=G(K)[p^\infty]$, by the \'etale
Kummer torsors and B1. Projection to $A$ and the preceding argument
put $d$ in $T(K)$. A torus has no nonidentity geometric $p$-power
 torsion points, since $z^{p^a}-1=(z-1)^{p^a}$ after splitting it.
Hence $d=0$. Push the point sections along the immersion $X\hookrightarrow G$
to obtain the assertion for $X$, as in B2. These comparisons use finite
\'etale covers throughout; a non-\'etale multiplication-by-$p$ torsor
has not been treated as an \'etale cover.
\end{proof}

\begin{result}{B7}{Refuted}{What geometrically prime-to-$p$ abelianization loses}
\label{res:B7}
For $K/\mathbb F_p$ finitely generated and $G/K$ semiabelian, the kernel
of the prime-to-$p$ Kummer map on $G(K)$ is exactly $G(K)[p^\infty]$.
Consequently prime-to-$p$ (**) is false even for hyperbolic
once-punctured elliptic curves. Without hyperbolicity, prime-to-$p$ (*)
is false too. We retain the full arithmetic quotient $G_K$ and replace
only the geometric kernel by its maximal prime-to-$p$ quotient.
This decides the explicit assertion of scan 35 (149, section 1,
manuscript p.~17; F7).
\end{result}
\begin{proof}
For $(n,p)=1$, $[n]:G\to G$ is a finite \'etale torsor under $G[n]$.
The comparison in B2 gives kernel $nG(K)$ at this level. The
prime-to-$p$ abelian geometric group of a semiabelian variety is
$\prod_{\ell\ne p}T_\ell G$: the multiplication covers yield these
quotients, as follows from the torus and abelian-variety Kummer
sequences. The compatible-division-point compactness argument of B2
therefore identifies the whole kernel with the intersection of these
finite-level kernels, namely $G(K)[p^\infty]$ by B1.

Take $E/\mathbb F_5:y^2=x^3+3x+2$. Its discriminant is nonzero and
\[
 E(\mathbb F_5)=\{O,(1,1),(1,4),(2,1),(2,4)\}.
\]
Thus this group has order five. The script \texttt{code/AB\_checks.py}
verifies the list and the orders by the group law. Put
$P=(1,1),Q=(1,4)$. For $Y=E\setminus\{O\}$ the reduced one-point
modulus has generalized Jacobian $E$, and the Abel map induces
\[
 \pi_1(Y_{\overline{\mathbb F}_5})^{\rm ab,(5')}
       \simeq\prod_{\ell\ne5}T_\ell E.
\]
Indeed the meridian module for a single puncture is zero. Hence
$P,Q$ have equal prime-to-five abelianized section classes on the
hyperbolic curve $Y$. For $X=E$ itself the geometric fundamental
group is abelian, so the same points disprove prime-to-five (*) in
the embedding generality of Theorem 3.

For a torus the displayed kernel is zero. Hence any variety immersed
in a torus has injective prime-to-$p$ (**) and (*). This includes
smooth genus-zero curves other than the projective and affine lines:
with a chosen rational point their generalized Jacobian is a torus
and the Abel map is an immersion. Passing to a finite extension to
split the punctures is legitimate: section equality restricts along
$G_{K'}\hookrightarrow G_K$, and inequality of points persists.
\end{proof}

\begin{result}{B8}{Proved}{Nonabelian recovery for punctured genus-one curves}
\label{res:B8}
For $K/\mathbb F_p$ finitely generated and $Y=C\setminus D$ a smooth
geometrically connected genus-one curve with nonempty reduced finite
boundary, prime-to-$p$ (*) is injective. Thus every hyperbolic genus-one
type is covered, including B7 where (**) fails. This proves part of
the covering argument on scan 36 (149, section 1, manuscript p.~17; F7).
\end{result}
\begin{proof}
Suppose $P,Q\in Y(K)$ give equal section classes. Use $P$ as origin
of $C=E$. Abelianization and B1 imply that $Q-P\in E(K)$ has
$p$-power order, say $p^a$. Suppose it is nonzero. Choose $(m,p)=1$
with $m^2>p^a$, and form
\[
 [m]:Y_m=E\setminus[m]^{-1}D\longrightarrow Y.
\]
Its geometric Galois group $E[m]$ has order $m^2$ prime to $p$.
Pass to one finite extension $L/K$ over which $E[m]$ and every
puncture are rational. The point $P'=0$ lies above $P$. Equal point
sections give a lift $Q'\in Y_m(L)$ with the same section class as
$P'$: rational lifts are fixed points of $G_L$ in the geometric fibre;
after conjugating the sections to be equal, the fixed fibre point for
$P$ is also fixed for $Q$. This proves the comparison rather than
assuming points lift over $K$.

The geometric prime-to-$p$ group of $Y_m$ is the corresponding open
subgroup of that of $Y$. Indeed, intersect a finite prime-to-$p$
Galois cover of $Y_m$ with its finitely many conjugates under $E[m]$.
The resulting Galois group over $Y$ is an extension of a subgroup of
a product of prime-to-$p$ groups by $E[m]$, and still has prime-to-$p$
order. Thus the induced topology is the full required topology.

Put $D_m=[m]^{-1}D$. Abelianization on $Y_m$ and B1 imply
$[Q'-P']\in J_{D_m}(L)[p^\infty]$. The kernel of $J_{D_m}\to E$
is a torus and has no geometric $p$-torsion. The order of this class
therefore equals the order of $Q'-P'$ in $E(L)$. Multiplication by
$m$ sends the latter to $Q-P$ and is invertible on $p$-primary groups;
the order is exactly $p^a$. The divisor description of the reduced
modulus Jacobian consequently gives a rational function $h$ with
\[
 \operatorname{div}(h)=p^a(Q'-P'),\qquad h|_{D_m}=1.
\]
The modulus condition requires the same nonzero value at all boundary
points, and rescaling gives the value $1$. The degree of $h$ is $p^a$,
yet $h-1$ vanishes at all $|D|m^2\geq m^2>p^a$ distinct points of
$D_m$, away from its poles. Equality of zero and pole degrees is a
contradiction. Hence $P=Q$.
\end{proof}

\begin{result}{B9}{Open}{The remaining prime-to-$p$ point question}
\label{res:B9}
We have not settled prime-to-$p$ (*) for every hyperbolic curve over
a field finitely generated over $\mathbb F_p$, nor the implication
on scan 36 that compatible $p$-primary differences on all
prime-to-$p$ covers force equality. The remaining range in this work
consists of proper and affine genus at least two with
nonzero $p$-primary Jacobian differences. Genus zero and punctured
genus one are settled by B7--B8. B14--B15 additionally separate the
two specific genus-two pairs over $\mathbb F_3$ and $\mathbb F_5$
investigated below. The reduction proved here sends every collision
over a finitely generated field to a collision over a finite field,
with the same geometric genus and number of punctures. Thus the
remaining universal question reduces exactly to smooth proper curves
of genus at least two over finite fields; already its universal
genus-two case remains undecided in this work.
This retains the unresolved portion
of F7 (149, section 1, manuscript p.~17, scans 35--36).
\end{result}
\begin{proof}[Verified partial result and obstruction]
\emph{The actual compactification comparison.}
Write $\Sigma=\{\ell:\ell\ne p\}$. For an affine hyperbolic
curve $U$ of genus at least two and its
smooth proper compactification $j:U\hookrightarrow C$, functoriality
of the geometric prime-to-$p$ quotient gives a continuous homomorphism
\[
 j_*:E_U^\Sigma\longrightarrow E_C^\Sigma
       \quad\text{over }\operatorname{id}_{G_K}.
\]
It sends the geometric kernel into the geometric kernel and sends
the point section at $R\in U(K)$ to its point section on $C$.
Thus an equality $s_Q=\operatorname{int}(\delta)s_P$ with
$\delta$ geometric gives the same equality on $C$, conjugated by
$j_*(\delta)$. The open immersion keeps $P,Q$ distinct.
Consequently every affine collision of genus at least two gives
a proper collision of the same genus. Proper curves are already
among the curves in question; together with B7--B8 this reduces the
universal problem to the proper genus-at-least-two case. The
finite-field comparison below then gives the claimed exact range.

Equal section classes do imply the covering condition: on every
finite extension where a chosen lift $P'$ is rational, the fixed-fibre
argument of B8 produces $Q'$ on any finite geometrically
prime-to-$p$ cover, and Kummer theory gives
$[Q'-P']\in J_{Y'}(L)[p^\infty]$.
The words ``cover of $X$'' on scan 36 must here refer to a cover of
$Y$, the target containing $P,Q$. This is evidence for editorial
correction, not a silently chosen reading. Here ``$p$-torsion'' means
$p$-primary torsion; no uniform exponent one has been established.

\emph{Reduction to finite fields, with the arithmetic maps.}
Let $U/K$ be a hyperbolic curve, where $K/\mathbb F_p$ is finitely
generated, and let $P\ne Q$ have equal geometrically prime-to-$p$
section classes. A smooth proper compactification with finite
\'etale boundary is part of the smooth-pair definition of a
hyperbolic curve. Even when only its geometric compactification is
initially given, its equations, boundary, and open immersion descend
from $\overline K$ to a finite extension $K'/K$. Replacing $K$ by
$K'$ preserves the distinct points and the equality of their section
classes. The actual field-base-change map is
\[
 E_{U_{K'}}^{\Sigma}\xrightarrow{\ \sim\ }
 E_U^{\Sigma}\times_{G_K}G_{K'},\qquad
 \Sigma=\{\ell:\ell\ne p\};
\]
for a separable extension this is the finite-\'etale base-change
description, and for a purely inseparable extension it follows from
the universal homeomorphism $U_{K'}\to U$. A finite extension is
handled by these two steps. In every group here the arithmetic
quotient is the full absolute Galois group.

Spread the compactification, its boundary, and the two sections over
a smooth integral affine scheme $S$ of finite type over a finite
field, with function field $K$. Shrinking $S$ gives a smooth proper
geometrically connected curve $C/S$, a finite \'etale divisor
$D\subset C$, and disjoint sections $P,Q$ of $\mathcal U=C\setminus D$.
These are finite-presentation conditions; the equalizer of the two
sections is closed and misses the generic point, so removing it
ensures disjointness. The genus and boundary degree are constant.
Choose a closed point $s\in S$, put
\[
 R=\mathcal O_{S,s}^{h},\qquad L=\operatorname{Frac}(R),
 \qquad k=\kappa(s).
\]
The field $k$ is finite. The ring $R$ is a regular henselian local
domain. The finite-\'etale equivalence for a henselian local ring
identifies $\pi_1(R)$ with $G_k$\cite[Tag 09ZS]{AB_Stacks}.
The homomorphism $r:G_L\to G_k$ induced by
$\operatorname{Spec}(L)\to\operatorname{Spec}(R)$ is surjective:
a connected finite \'etale cover of the normal integral scheme
$\operatorname{Spec}(R)$ is normal and integral, and hence remains
connected over $L$.

Here are the required relative groups, without replacing the
arithmetic quotient by a prime-to-$p$ quotient. Set
\[
 K_R=\ker\bigl(\pi_1(\mathcal U_R)\to\pi_1(R)\bigr),\qquad
 N_R=\ker(K_R\to K_R^\Sigma),\qquad
 E_R=\pi_1(\mathcal U_R)/N_R.
\]
This is exactly the construction of SGA1 XIII, 4.0.
Proposition 4.3 gives the homotopy sequence; Examples 4.4 verify
its hypotheses for a smooth proper family with relative
normal-crossings boundary and a section. Section 4.5.3 identifies
the relative pro-$\Sigma$ kernel after every connected base change,
in the source's words ``canoniquement isomorphe \`a l'image inverse''
\cite[Expos\'e XIII, 4.0, Proposition 4.3, Examples 4.4, 4.5.3]{SGA1}.
Our finite \'etale boundary is such a relative divisor; the section
$P$ supplies the required section. With compatible base points and
paths supplied by $P$, the closed- and generic-fibre inclusions give
\[
 \begin{gathered}
 1\longrightarrow\Delta_s^\Sigma\longrightarrow E_R
   \longrightarrow G_k\longrightarrow1,\qquad
 j_s:E_{U_s}^{\Sigma}\xrightarrow{\ \sim\ }E_R,\\
 \operatorname{sp}=j_s^{-1}j_L:E_{U_L}^{\Sigma}
   \longrightarrow E_{U_s}^{\Sigma},\\
 \theta:E_{U_L}^{\Sigma}\xrightarrow{\ \sim\ }
 E_{U_s}^{\Sigma}\times_{G_k}G_L,
 \qquad e\longmapsto\bigl(\operatorname{sp}(e),\operatorname{pr}(e)\bigr).
 \end{gathered}
\]
Indeed, $j_s$ is an isomorphism on both the geometric kernel and
the arithmetic quotient. The map $j_L$ is an isomorphism on the
geometric kernel by 4.5.3; consequently $\theta$ is an isomorphism
on its kernel and quotient, proving the asserted isomorphism of
profinite groups. These descriptions also prove surjectivity of
$\operatorname{sp}$.

Naturality for the two actual sections of $\mathcal U_R/R$ gives
\[
 \operatorname{sp}\circ s_{P_L}=s_{P_s}\circ r,
 \qquad
 \operatorname{sp}\circ s_{Q_L}=s_{Q_s}\circ r.
\]
The equality of the original section classes restricts to $L$.
Thus, after compatible geometric choices,
$s_{Q_L}=\operatorname{int}(\delta)s_{P_L}$ for
$\delta\in\Delta_L^\Sigma$. Applying $\operatorname{sp}$ and using
surjectivity of $r$ gives
\[
 s_{Q_s}=\operatorname{int}\bigl(\operatorname{sp}(\delta)\bigr)s_{P_s}.
\]
The conjugator belongs to $\Delta_s^\Sigma$, and $P_s\ne Q_s$ by
construction. We have produced a finite-field collision of the
same type. Conversely a finite field is itself finitely generated
over $\mathbb F_p$. This proves equivalence of the two universal
injectivity assertions, separately for each geometric curve type.

The opened finite-field result of Sa\"idi--Tamagawa is precisely
\cite[Proposition 1.8(vi), Corollary 1.11(iii)]{AB_SaidiTamagawa2009}:
for a hyperbolic curve and a cofinite retained set of primes, the
closed-point decomposition map has a finite exceptional set. The
introduction, p.~138, explicitly says ``almost injective''. The 2017
refinement retains the exceptional set\cite[Definitions 1.1--1.4,
Theorem 1.9]{AB_SaidiTamagawa2017}. Both texts were opened.
Equality of point sections implies conjugacy of decomposition groups,
so a possible coincidence lies in that finite set. Neither statement
makes it empty. The Isom theorem reconstructing the curve does not
supply the missing implication about this point map.

We can prove exactly why the finite exceptional set does not vanish
merely by passing to covers. Use the notation of
\cite[Remark 1.2, Definition 1.5, Proposition 1.8]{AB_SaidiTamagawa2009}.
For a hyperbolic $U$ over a finite field $k$, write $\Pi_U$ for its
geometrically prime-to-$p$ arithmetic group and $\widetilde U$ for
the associated universal pro-cover. For a chosen pro-point $x$,
let $D_x\subset\Pi_U$ be its decomposition subgroup. The exceptional
pro-point set $\widetilde E_U$ consists exactly of the $x$ for which
there is $y\ne x$ with $D_x,D_y$ commensurate. For an open subgroup
$H\subset\Pi_U$, with corresponding curve $V$, the same pro-cover
serves $V$, and its decomposition subgroup is $D_x\cap H$.
Intersection with $H$ preserves and reflects commensurability:
each $D_x\cap H$ is open in $D_x$, and a common finite-index
intersection in $H$ is also one in $\Pi_U$. Hence
\[
 \widetilde E_V=\widetilde E_U,\qquad
 E_V=\operatorname{preimage}(E_U)\quad\hbox{in }V.
\]
The second equality uses the $\Pi_U$-stability in Proposition1.8(v)
and the orbit descriptions $E_U=\widetilde E_U/\Pi_U$,
$E_V=\widetilde E_U/H$. For a geometrically connected degree-$d$
cover, on geometric point sets it gives
$|E_V(\overline k)|=d|E_U(\overline k)|$. Finiteness for each cover
therefore supplies no shrinking argument or uniform bound.

Moreover $E_U$ is empty exactly when the point-section map is
injective over every finite constant extension of $k$. To prove the
nontrivial direction, take distinct exceptional pro-points $x,y$
and put $Z=D_x\cap D_y$. The cusps are excluded by Proposition1.8(ii),
so the projections of $D_x,D_y$ to $G_k$ are injective. The subgroup
$G_0=\operatorname{pr}(Z)$ is open in $G_k$; on restricting to
$\operatorname{pr}^{-1}(G_0)$, both decomposition subgroups become
exactly $Z$, by injectivity of these projections. Over the finite
constant extension fixed by $G_0$, the two images are rational
points with the same section. They are distinct, since
Proposition1.8(iii) says commensurate decomposition groups cannot
occur at distinct pro-points above the same geometric point.
Conversely equal point-section classes can be represented by equal
actual decomposition groups after choosing the corresponding
pro-lifts, and give an exceptional pair. Thus an actually nonempty
exceptional set would give a counterexample after a finite extension;
the cited results do not assert that such a set is nonempty.

The first concrete proper test we attacked was the smooth model
$y^3-y=x+x^{-1}$ over $\mathbb F_3$. It has genus two and two rational
points $P,Q$ over $0,\infty$, with $\operatorname{div}(x)=3P-3Q$.
At each pole, $1/y$ is a uniformizer and
$v(1/(y+c)-1/y)=2$ for $c\in\mathbb F_3^*$; the two different
exponents are four, giving $2g-2=3(-2)+8=2$.
B14 now separates these sections by an explicit \'etale double
cover over $\mathbb F_{27}$: no lift of $Q$ has $3$-primary Jacobian
difference from the chosen lift of $P$. This test is resolved, and
is not retained as an undecided example or as a counterexample.

The second explicit genus-two test is the smooth
projective model over $\mathbb F_5$ of
\[
 C:\quad y^2=x^5+(x^2+1)^2,\qquad P=(0,1),\quad Q=\infty.
\]
The polynomial on the right is squarefree: its derivative is
$4x(x^2+1)$, whose zeros are not zeros of the polynomial. It has
degree five, so the compactification is a smooth genus-two curve
with its unique rational point $Q$ at infinity. At $P$, $x$ is a
uniformizer and
\[
 (y-x^2-1)(y+x^2+1)=x^5
\]
gives $\operatorname{div}(y-x^2-1)=5P-5Q$: the second factor is
nonzero at $P$, and $y-x^2-1$ has pole order five at infinity.
Thus $[P-Q]$ has exact order five, since a function with divisor
$P-Q$ would have degree one. The squarefreeness and finite point
calculation are included in \texttt{code/AB\_checks.py}.
B15 separates their prime-to-five section classes. Explicitly,
write $x^5+(x^2+1)^2=(x-1)H(x)$, where
$H(x)=x^4+2x^3+2x^2-x-1$. The \'etale double cover obtained by
adjoining $v^2=x-1$ has the elliptic quotient
$E:w^2=H(x)$, with $w=y/v$ and $|E(\mathbb F_5)|=8$.
A chosen lift of $P$ maps to $(0,3)$, while the two lifts of $Q$
map to the two infinity points. Both resulting elliptic differences
are nonzero of $2$-power order, contradicting the necessary
$5$-primary condition. The complete cover and quotient verification
is B15. Thus neither explicit pair is left undecided. The remaining
statement is the universal higher-genus finite-field assertion,
including genus two with nonzero $p$-primary Jacobian differences.

\emph{An exact finite test in odd-characteristic genus two.}
The two constructions also give a verified next test, with its
remaining obligation explicit. The canonical linear system of a
genus-two curve has dimension one and degree two and is base-point
free: a base point $R$ would contradict Riemann--Roch, which gives
$h^0(K_C-R)=h^0(R)=1$. It defines a separable degree-two map to
$\mathbb P^1$ when $p$ is odd. After a finite constant extension
splitting its six branch points, a projective change of coordinate
puts infinity outside that finite set. In these actual coordinates
write the curve as $z^2=A_2(x)B_4(x)$ with six distinct
branch points, choosing any pair for $A_2$. There are fifteen
choices of the pair. The curve
$D:v^2=A_2(x),\ w^2=B_4(x)$ is a connected \emph{\'etale}
double cover, with elliptic quotient $E:w^2=B_4(x)$; the proof is
the same even-divisor and conic-quotient calculation as B14, with
the displayed branch polynomials in place of its particular ones.
The two involutions have hyperelliptic product, so the exact norm
identity $f^*f_*+g^*g_*=[2]$ holds on $J_D$.
Make a chosen lift $P'$ and both lifts $Q'$ rational over one
finite field. For $d=[Q'-P']$, the two norms are
$[Q-P]$ and $[g(Q')-g(P')]$. A non-$p$-primary elliptic norm
already separates that lift, by the finite Lang construction of
B14. When both norms have respective $p$-power orders $p^a,p^b$,
put $N=p^{\max(a,b)}$. The norm identity gives $Nd\in J_D[2]$.
Since $p$ is odd, $d$ is $p$-primary exactly when $Nd=0$:
in the forward direction $Nd$ is both $p$-primary and killed by
two, and in the reverse direction $N$ is a power of $p$.
The last equality is an exact divisor-class test; it is equivalent
to a nonzero section of the degree-zero line bundle
$\mathcal O_D(NP'-NQ')$, and hence is decidable by a finite
Riemann--Roch computation. This identifies the residual obstruction
in each of the fifteen first double-cover tests. We have not proved
that one of those tests always separates a distinct pair, or that
successive deeper covers eliminate every pair for which they fail.
The examples B14--B15 verify success for their two specified pairs;
no universal success is inferred from them.

\emph{Compatibility of the fifteen lifts and their exact joint kernel.}
There is a further proved reduction of these fifteen tests. Work in
odd characteristic, put $J=J_C$, and use the Abel map
$j:C\to J$, $R\mapsto[R-P]$. Its map on geometric fundamental
groups identifies the abelianization with that of $J$
\cite[III, Proposition 9.1]{AB_MilneAV}. This opened proposition
works over a separably closed field; only its prime-to-characteristic
case $n=2$, proved there by Lemmas 9.2--9.3, is used here.
Consequently
\[
 Z=C\times_{j,J,[2]}J,\qquad \pi:Z\longrightarrow C,
 \qquad G=J[2](\overline k)\simeq\mathbb F_2^4
\]
is a connected \'etale Galois cover of degree sixteen and genus
seventeen. Its fifteen index-two quotients are precisely the
fifteen double covers just considered. Here is the exact
identification with the branch-pair list. Over $\overline k$,
Kummer theory identifies connected double covers with the nonzero
classes in $J_C[2]$\cite[III, Lemma 9.2]{AB_MilneAV}.
The fifteen classes $[W_i-W_j]$ of distinct Weierstrass points
are nonzero and distinct as unordered pairs. Nonzero follows
because a function with divisor $W_i-W_j$ would have degree one.
They are killed by two since $2W_i$ and $2W_j$ are fibres of
the hyperelliptic map. For distinctness, equality of two such
classes gives $W_i+W_l\sim W_k+W_j$. If these effective
degree-two divisors differ, Riemann--Roch makes them canonical
fibres. A canonical fibre supported on Weierstrass points is
$2W$, forcing $i=l$ and $k=j$, hence the same unordered pair.
Equality of the effective divisors themselves gives that pair
as well. Since $|J_C[2]|=16$, the list is exhaustive.
In the degree-six model with infinity unramified, the divisor
of $(x-a_i)(x-a_j)$ is
$2W_i+2W_j-2\infty_+-2\infty_-$; its Kummer class is
$[W_i+W_j-\infty_+-\infty_-]=[W_i-W_j]$.
Thus the displayed square-root construction gives exactly these
quotients. All the groups, covers and
points below are defined over one finite extension of the original
finite field. This follows by descending their finite algebraic
data from its algebraic closure.

For a hyperplane $H\subset G$, put
\[
 B_H=J/H,\qquad \rho_H:J\to B_H,\qquad
 [2]=\lambda_H\rho_H,\qquad
 D_H=C\times_{j,J,\lambda_H}B_H.
\]
Here $\lambda_H:B_H\to J$ is an isogeny of degree two,
$f_H:D_H\to C$ is the corresponding double cover, and
\[
 r_H:Z\longrightarrow D_H,
 \qquad (R,b)\longmapsto(R,\rho_H(b))
\]
is the quotient by $H$, of degree eight. Choose
$\widetilde P=(P,0)$ and let $P_H=r_H(\widetilde P)$.
Changing the chosen lift of $P$ in an individual double cover
applies its deck automorphism to the two tested differences and
does not change whether a $p$-primary lift exists.

Let $c=[Q-P]$ be $p$-primary. Multiplication by two is an
automorphism of the $p$-primary torsion group, so there is a unique
$p$-primary $b$ with $2b=c$; explicitly $b=mc$ for an integer
$m$ inverse to two modulo the order of $c$. Put
$\widetilde Q=(Q,b)\in Z$. Every successful lift in the fifteen
tests must be the image of this single point. Indeed the second
projection $a_H:D_H\to B_H$ sends $P_H$ to zero and hence induces
a homomorphism $J_{D_H}\to B_H$. If
$[Q_H'-P_H]$ is $p$-primary, then $a_H(Q_H')$ is $p$-primary
and maps to $c$ under $\lambda_H$. An isogeny of degree two
has exactly one $p$-primary point over a given $p$-primary point:
its kernel has exponent two, and $\rho_H(b)$ supplies that
preimage. Thus $a_H(Q_H')=\rho_H(b)$. In the displayed fibre
product, the coordinate pair $(Q,a_H(Q_H'))$ determines $Q_H'$
itself. Therefore
\[
       Q_H'=r_H(\widetilde Q).
\]
In particular each test has at most one successful lift, and
success in all fifteen tests makes their choices compatible.

Write $\delta=[\widetilde Q-\widetilde P]\in J_Z$.
The exact equivalence just proved is that all fifteen tests
admit a $p$-primary lift precisely when every
$(r_H)_*\delta$ is $p$-primary. To see what those norm data
determine, write $N_A=\sum_{a\in A}a_*$ for a subgroup
$A\subset G$. In the integral group ring the hyperplane
incidence identity is
\[
 \sum_{\substack{H\subset G\\ \lvert G/H\rvert=2}}N_H
       =8\operatorname{id}+7N_G.
\]
The identity element belongs to fifteen hyperplanes and every
other element to seven. Pullback and norm therefore give the
following exact equality of endomorphisms of $J_Z$:
\[
 [8]=\sum_H r_H^*(r_H)_*-7\pi^*\pi_*.
\]
Let
\[
        K_2=\bigcap_H\ker\bigl((r_H)_*:J_Z\to J_{D_H}\bigr).
\]
Since $\pi_*= (f_H)_*(r_H)_*$, this identity proves
$K_2\subset J_Z[8]$. In particular it is finite and
$2$-primary. When all fifteen tests succeed, choose one power
$p^a$ killing their fifteen norm classes. For the actual
point-divisor class above we obtain
\[
          p^a\delta\in K_2,\qquad 8p^a\delta=0.
\]
Thus the joint tests leave only a $2$-primary part killed by
eight in this degree-sixteen cover.

This joint kernel really is nonzero for every such genus-two
curve. The character map
\[
 G^D=\operatorname{Hom}(G,\mathbb G_m)
 \longrightarrow\operatorname{Pic}^0(C),\qquad
 \chi\longmapsto Z\times^G\mathbb A^1_\chi
\]
identifies $G^D$ with $\ker(\pi^*:J_C\to J_Z)$.
For a line bundle trivialized upstairs, the descent multipliers
belong to $\Gamma(Z,\mathcal O_Z^*)=\overline k^*$ and
give exactly such a character. Conversely a character gives the
displayed associated line bundle. This map is injective: a
trivial descended bundle would give a nonzero global function
upstairs transforming by that character, whereas all such units
are constant. Every character has order dividing two and its
associated bundle has degree zero. There are sixteen characters;
$J_C[2]$ also has sixteen geometric points. As these are finite
\'etale groups in odd characteristic, this proves the exact equality
\[
                  \ker\pi^*=J_C[2].
\]
It follows that the actual image
\[
 \pi^*J_C[8]\simeq J_C[8]/J_C[2]
                   \simeq(\mathbb Z/4)^4
\]
is a nonzero subgroup of $J_Z$. For every $a\in J_C[8]$,
the factorization $\pi=f_Hr_H$ gives
\[
 (r_H)_*\pi^*a
   =(r_H)_*r_H^*f_H^*a=8f_H^*a=0.
\]
Hence $\pi^*J_C[8]\subset K_2$. In particular an $a$ of
exact order eight has $\pi^*a$ of exact order four: its double
cannot vanish, since $2a$ has order four and is not in
$\ker\pi^*=J_C[2]$.

The remaining point-divisor implication is precise. The fifteen
conditions place the specific class
$p^a[\widetilde Q-\widetilde P]$ in this nonzero finite kernel;
they do not, by the norm identities alone, determine its value.
We have not proved that the geometry of these two lifted points
forces the residual class to vanish, or that success of all fifteen
tests forces $P=Q$. Nor has a nonzero class constructed above been
identified with such a point difference satisfying all the tests.
The hyperplane identity is checked independently by the
$15$-by-$16$ incidence matrix in \texttt{code/AB\_checks.py}.

The genus-one argument succeeds because the compactification of each
$[m]^{-1}Y$ remains the same elliptic curve; the $p$-power order is
fixed while the number of punctures grows. Higher-genus covers
introduce additional Jacobian factors, and our argument does not
bound the $p$-primary orders there. Proper covers also have no modulus
points on which to force $h=1$. These are the precise remaining
obstructions to the proposed covering argument.
\end{proof}

\begin{result}{B10}{Proved}{Marked field faithfulness in every characteristic}
\label{res:B10}
For any field $k$, the contravariant functor from finitely generated
field extensions of $k$ to their profinite \'etale fundamental groupoids
over $\Pi_k$ is faithful, where an isomorphism between two functors
must project to the identity natural transformation on $\Pi_k$.
In characteristic $p>0$ it remains faithful on replacing each geometric
kernel by its maximal prime-to-$p$ quotient and retaining the arithmetic
base groupoid. Thus the characteristic-$p$ margin to Theorem 2 is true
in this precise marked sense, in particular for finitely generated
extensions of $\mathbb F_p$ (149, section 1, manuscript pp.~10--11,
17, scans 22--23 and 36; F7).
\end{result}
\begin{proof}
Let $f,g:F\rightrightarrows L$ be $k$-embeddings and suppose their
induced groupoid functors are isomorphic over $\Pi_k$ with the stated
identity condition. Base change the groupoids to a geometric point of
$\Pi_k$. The components correspond to embeddings into $\overline k$
of the finite relative separable constant fields. The given natural
isomorphism identifies the component maps. On each corresponding
component it identifies the induced geometric homomorphisms up to
inner conjugation. Purely inseparable extensions do not change the
category of finite \'etale covers: a finite separable extension and
its embeddings extend uniquely after such a base change, and descend
by taking the separable elements. We may therefore pass to
$\overline k$ and work with the field components of $F\overline k$
and $L\overline k$. The extended embeddings fix $\overline k$;
restriction back to $F$ will give the original equality.

For an integer $n$ invertible in $k$, Kummer theory gives
\[
 H^1(G_{F\overline k},\mu_n)
       =(F\overline k)^*/(F\overline k)^{*n}.
\]
The coefficients are constant, and these degree-one classes factor
through the geometric prime-to-$p$ quotient when $p>0$. Inner
conjugation acts trivially on this $H^1$. Thus for every $a\ne0$
we have $f(a)/g(a)\in(L\overline k)^{*n}$ for every such $n$.
On a normal proper model of the finitely generated field
$L\overline k/\overline k$, every divisorial valuation of this ratio
is divisible by arbitrarily large integers prime to the characteristic,
and hence is zero. A rational function with zero divisor on a normal
proper integral variety is an invertible global function, and is
constant because the base is algebraically closed. Consequently
$f(a)=c_a g(a)$ for some $c_a\in\overline k^*$.

For nonconstant $a$ apply this also to $a-1$. Writing $c=c_a$ and
$d=c_{a-1}$ gives
\[
       (c-d)g(a)=1-d.
\]
The element $g(a)$ is nonconstant, so $c=d=1$. Constants are fixed
already. Therefore the two extended field maps agree, and hence
$f=g$. This proof constructs the comparison through the Kummer
classes; it does not rely on an unproved reduction to a product of
curves.

The arithmetic marking is material. On $\mathbb F_p(t)$, Frobenius
$t\mapsto t^p$ acts on the Kummer class of $t$ by multiplication by
$p$. Choose $n$ prime to $p$ and not dividing $p-1$; the $t$-adic
valuation shows $[t^p]\ne[t]$ modulo $n$th powers. Thus the marked
functors differ. The usual natural equivalence on the unmarked
\'etale site coming from absolute Frobenius has component on the
chosen algebraic constant field equal to arithmetic Frobenius, not
to the identity. It consequently is not an allowed relative natural
transformation in Theorem 2. This also explains the multiplication-by-$p$
action on $\pi_1(\mathbb G_{m,\overline{\mathbb F}_p})^{(p')}$:
it is conjugation by an arithmetic Frobenius lift, whose image in
$G_{\mathbb F_p}$ is nontrivial.
\end{proof}
